\begin{question}
\noindent A satellite dish engineer models the cross-section of a circular antenna dish as a circle with centre $O$. Points $A$, $B$ and $C$ lie on the circumference, where $AC$ is a diameter of the circle. A support strut $CD$ is attached so that it is tangent to the circle at point $C$.

\begin{center}
\begin{tikzpicture}[scale=2.2]
    % Circle centre O at origin, radius 1
    \coordinate (O) at (0,0);
    \coordinate (A) at (180:1);
    \coordinate (C) at (0:1);
    \coordinate (B) at (130:1);

    \draw[name path=circ] (O) circle (1);
    \fill[black] (O) circle (0.4pt);

    % Diameter AC
    \draw[thick] (A) -- (C);

    % Chords AB and BC
    \draw[thick] (A) -- (B);
    \draw[thick] (B) -- (C);

    % Tangent at C, perpendicular to AC (i.e. vertical line)
    \draw[thick] (C) -- ++(90:1) coordinate (Dtop);
    \draw[thick] (C) -- ++(-90:1) coordinate (Dbot);

    % Labels
    \node[left] at (A) {$A$};
    \node[above right] at (B) {$B$};
    \node[right] at (C) {$C$};
    \node at (O) [below] {$O$};
    \node[above right] at (Dtop) {$D$};

    % angle at B (angle in semicircle)
    \pic [draw, angle eccentricity=1.4, angle radius=0.45cm, "$90^\circ$"] {angle = C--B--A};

    % angle BAC = 35 degrees marked at A
    \pic [draw, angle eccentricity=1.5, angle radius=0.35cm, "$35^\circ$"] {angle = B--A--C};

    % right angle at C between AC and tangent CD
    \tkzMarkRightAngle[size=0.25](A,C,Dtop);

\end{tikzpicture}
\end{center}

\noindent Angle $BAC = 35^\circ$.

\begin{enumerate}
    \item[(a)] Show that angle $ABC = 90^\circ$, giving a reason for your answer.

    \vspace{3cm}
    \answermarks{2}

    \item[(b)] Work out the size of angle $BCD$.\\
    Give a reason for each stage of your working.

    \vspace{4cm}
    \answerequnit{BCD}{$^\circ$}{2}
\end{enumerate}
\end{question}
